\hypertarget{choosing-an-r-discovery-and-effect-set-boundary}{%
\section{Choosing an R discovery and effect-set
boundary}\label{choosing-an-r-discovery-and-effect-set-boundary}}

\label{exec13:pcalg-reuse}

The preferred R route is a restricted integration of \texttt{pcalg} for
Gaussian observational discovery and singleton linear effects. Its
attraction is the connection between a partially directed graph and a
multiset of effects: the product can preserve uncertainty about
direction instead of forcing a unique business answer. The proposal
selects this pinned path for local prototype implementation under the
strict contract developed below. Runtime qualification remains open
until the environment and conversion boundary have been exercised;
selection does not claim that the proposed R adapter already works.

The selected source is \texttt{pcalg} version 2.7-12, dated 12 September
2024, at CRAN mirror revision
\texttt{930fe476875ddc00528dbaacb6f48f2791ab12e3}. The bounded history
query selects the most recent mirror commit before the 7 October 2026
cutoff. The reference manual was generated in May 2026; this later
document date does not change the release into a 2026 implementation.
Ten inspected source files have byte identities matching the selected
Git tree. The mirror is explicitly a read-only copy of the CRAN package,
and its commit history establishes package identity and release
activity, rather than upstream support quality
\citep{exec13_pcalg,exec13_pcalg_reuse_primary}.

PC, its stable and conservative variants, FCI, GIES and IDA are
operations or families within this package, not independent
implementations whose agreement would multiply evidence. The initial
integration selects only the PC/IDA path. The presence of
latent-variable and interventional algorithms in the same package does
not make their model classes interchangeable. A later FCI route would
require a separate partial ancestral graph contract, and a later GIES
route would require intervention metadata and score assumptions. A
successful installation cannot settle either scientific extension.

\hypertarget{what-an-effect-multiset-would-mean-for-a-bank}{%
\subsection{What an effect multiset would mean for a
bank}\label{what-an-effect-multiset-would-mean-for-a-bank}}

Let the observed vector contain an offered rate change \(X\), an
intermediate customer response \(M\), and a monetary outcome \(Y\). The
initial model is a linear Gaussian structural system with independently
sampled units, no hidden common causes and a distribution faithful to
its causal DAG. The sampling unit could be a customer in a stipulated
independent-unit experiment; repeated customer months cannot be
relabeled as independent rows. The graph search estimates conditional
independence structure. Its significance level is a level for individual
tests, not a confidence level for the entire graph or for the effect of
a future bank action \citep{exec13_pcalg}.

For a compatible DAG \(G\), the singleton total effect is the derivative
\[
 \theta_G=\frac{\partial}{\partial x}
     \mathbb E_G[Y\mid\operatorname{do}(X=x)].
\] In the stated linear model it is constant in \(x\). When \(Y\) is not
a parent of \(X\), it can be calculated by regressing \(Y\) on \(X\) and
the parents of \(X\) in \(G\). If \(Y\) is a parent, its total effect
from \(X\) is zero. The global IDA construction considers compatible
DAGs; the local construction uses admissible local parent
configurations. They need not attach the same multiplicity to a
numerical value. Neither construction turns a frequency in this multiset
into a posterior probability of a graph
\citep{exec13_pcalg,exec13_pcalg_reuse_primary}.

The distinction is visible without fitting a package. Consider the exact
covariance matrix, ordered as \((X,M,Y)\), \[
 \Sigma=\begin{pmatrix}
  1&0.6&0.48\\
  0.6&1.36&1.088\\
  0.48&1.088&1.8704
 \end{pmatrix}.
\] The partial covariance of \(X\) and \(Y\) given \(M\) is
\(0.48-0.6(1.088)/1.36=0\). The three DAGs \(X\to M\to Y\),
\(X\leftarrow M\to Y\), and \(X\leftarrow M\leftarrow Y\) can each
generate this same Gaussian law with positive independent error
variances. Exact rational reconstruction of all three systems gives
total effects \(0.48,0,0\). Global enumeration therefore has multiset
\([0,0,0.48]\); the two local parent configurations of \(X\) give
\([0,0.48]\). Averaging the first would produce \(0.16\), and averaging
the second would produce \(0.24\). Neither average is identified as the
bank's intervention effect. These are exact project constructions and
independent references, not results from executing \texttt{pcalg}
\citep{exec13_pcalg_reuse_examples}.

Suppose \(X\) is measured in percentage points and \(Y\) in thousands of
currency units. The nonzero derivative becomes \(4.8\) currency units
per basis point, so a 25-basis-point action gives 120 currency units per
modeled unit in the forward-chain system. The other systems permit zero.
In general, for \(X'=c_X X\) and \(Y'=c_Y Y\), \[
 \Sigma'=D\Sigma D^{\mathsf T},\qquad
 \theta'_G=(c_Y/c_X)\theta_G.
\] Both the covariance transformation and the output units must survive
an R process boundary. A correlation matrix alone is insufficient to
restore monetary units. Further, naming \(X\) an offered rate is not
evidence for its causal position. Justified timing or an intervention
can exclude competing graphs; vocabulary alone cannot. Sampling
uncertainty remains a separate obligation even when a graph class has
been correctly represented.

\hypertarget{matrix-conventions-change-causal-meaning}{%
\subsection{Matrix conventions change causal
meaning}\label{matrix-conventions-change-causal-meaning}}

An adjacency matrix is not a complete interface specification. In this
package, the documented \texttt{amat.cpdag} convention refers to the row
endpoint, whereas \texttt{amat.pag} refers to the column endpoint. The
selected \texttt{pcAlgo} coercion calls \texttt{wgtMatrix}, whose
default transposes the underlying graph matrix. Thus a consumer that
changes the accessor while keeping its old decoder can reverse every
directed edge \citep{exec13_pcalg,exec13_pcalg_reuse_primary}.

For distinct named nodes \(i\) and \(j\), write the ordered pair
\((A_{ij},A_{ji})\). The following conventions must remain separate.

\begin{longtable}[]{@{}llll@{}}
\toprule
\begin{minipage}[b]{0.22\columnwidth}\raggedright
Representation\strut
\end{minipage} & \begin{minipage}[b]{0.22\columnwidth}\raggedright
Pair for \(i\to j\)\strut
\end{minipage} & \begin{minipage}[b]{0.22\columnwidth}\raggedright
Pair for an undirected edge\strut
\end{minipage} & \begin{minipage}[b]{0.22\columnwidth}\raggedright
Interpretation requirement\strut
\end{minipage}\tabularnewline
\midrule
\endhead
\begin{minipage}[t]{0.22\columnwidth}\raggedright
\texttt{amat.cpdag} from the \texttt{pcAlgo} matrix coercion\strut
\end{minipage} & \begin{minipage}[t]{0.22\columnwidth}\raggedright
\((0,1)\)\strut
\end{minipage} & \begin{minipage}[t]{0.22\columnwidth}\raggedright
\((1,1)\)\strut
\end{minipage} & \begin{minipage}[t]{0.22\columnwidth}\raggedright
Row-endpoint convention; retain the exact accessor.\strut
\end{minipage}\tabularnewline
\begin{minipage}[t]{0.22\columnwidth}\raggedright
Matrix of the underlying \texttt{graphNEL} object\strut
\end{minipage} & \begin{minipage}[t]{0.22\columnwidth}\raggedright
\((1,0)\)\strut
\end{minipage} & \begin{minipage}[t]{0.22\columnwidth}\raggedright
\((1,1)\)\strut
\end{minipage} & \begin{minipage}[t]{0.22\columnwidth}\raggedright
Row denotes source; verify the installed graph conversion.\strut
\end{minipage}\tabularnewline
\begin{minipage}[t]{0.22\columnwidth}\raggedright
\texttt{amat.pag}\strut
\end{minipage} & \begin{minipage}[t]{0.22\columnwidth}\raggedright
\((2,3)\)\strut
\end{minipage} & \begin{minipage}[t]{0.22\columnwidth}\raggedright
\((3,3)\)\strut
\end{minipage} & \begin{minipage}[t]{0.22\columnwidth}\raggedright
Column-endpoint marks; outside the initial effect path.\strut
\end{minipage}\tabularnewline
\bottomrule
\end{longtable}

A symmetric example cannot detect a transpose. The proposed acceptance
fixtures therefore include a named directed chain and a named collider,
each under every permutation of its three columns. An independently
written decoder was checked on 54 representation/fixture/permutation
combinations, and 36 deliberate transposes of directed fixtures changed
the expected answer. These checks establish the small reference
examples. They do not establish that an installed R serializer preserves
them; that is the later adapter test
\citep{exec13_pcalg_reuse_examples}.

The response should retain the original representation, node order,
exact accessor and package revisions alongside named edges. A display
may render the graph more conveniently, but an unresolved edge must stay
unresolved. The importer must also retain the origin of exceptional
marks. A pair of twos can encode an orientation conflict in PC and a
bidirected edge in a PAG; equal integers do not establish equal
scientific meaning.

\hypertarget{failures-must-be-caught-before-they-are-hidden}{%
\subsection{Failures must be caught before they are
hidden}\label{failures-must-be-caught-before-they-are-hidden}}

The selected source makes three failure boundaries especially
consequential. First, with \texttt{solve.confl=false}, relaxed PC can
overwrite an earlier edge direction when another collider or orientation
rule conflicts with it. With \texttt{solve.confl=true}, it can preserve
conflicting directions using reciprocal entries of two. The proposed
profile therefore uses conservative PC, stable skeleton search, relaxed
orientation and conflict preservation. Such a result may still fail to
be a valid CPDAG. The discovery result can retain this diagnostic, while
the effect calculation must refuse it. The conflict is not evidence of a
latent common cause \citep{exec13_pcalg_reuse_primary}.

Second, a finite conditional-independence p value need not show that
every preceding numerical operation was valid. The inspected
\texttt{pcorOrder} uses a pseudoinverse for conditioning sets of size at
least two, replaces a missing partial correlation by zero, and clips its
result. \texttt{zStat} likewise replaces a missing statistic by zero;
\texttt{gaussCItest} then evaluates a two-sided Gaussian tail
probability. Disabling \texttt{NAdelete} acts later in the skeleton
routine and cannot recover a failure already replaced inside these
functions \citep{exec13_pcalg_reuse_primary}.

For example, in the invalid correlation matrix with every entry equal to
one, the one-variable conditioning formula gives \[
 r_{12\mid3}
   =\frac{1-1\cdot1}{\sqrt{(1-1^2)(1-1^2)}}=\frac00.
\] Replacing this undefined value by zero yields a finite p value of
one. An independently written base-R scalar trace reproduces this
arithmetic; it also verifies that positive-definite factorization
rejects the input. A separate invalid-degrees-of-freedom example reaches
the same finite p value through the statistic fallback. These are
intentionally invalid inputs, already excluded by a correctly enforced
preflight. They demonstrate why checking only the final p value is
insufficient, without establishing a failure for valid, well-conditioned
package inputs \citep{exec13_pcalg_reuse_examples}.

The proposed initial path instead supplies a strict Gaussian test
callback through PC's existing independence-test interface. It uses
checked positive definite submatrices and factorization solves,
preserves degrees of freedom, and refuses rank loss, nonpositive
residual variances, nonfinite arithmetic or unacceptable conditioning.
It does not silently substitute a pseudoinverse, clip a failed
calculation or choose a new regularization parameter. Numerical
tolerances must be fixed before adapter qualification. This callback is
a proposed project component with its own comparison obligation, rather
than a description of unmodified \texttt{gaussCItest}. The compiled fast
skeleton path is outside the initial profile so that it cannot bypass
this callback.

Third, the inspected IDA code replaces all nonzero graph weights by one
before checking graph validity. If validity fails, it emits a message
without an immediate stop. A conflict pair \((2,2)\) can therefore lose
its distinction from an undirected pair \((1,1)\) before that check. The
proposed adapter must inspect the untouched graph, reject conflict
codes, and enforce a strict CPDAG/extension gate before calling IDA. It
must not treat a warning stream as a structured failure response, or
silently orient an edge to make an invalid request calculable
\citep{exec13_pcalg_reuse_primary}.

\hypertarget{reuse-choice-operating-cost-and-the-next-experiment}{%
\subsection{Reuse choice, operating cost and the next
experiment}\label{reuse-choice-operating-cost-and-the-next-experiment}}

An isolated local R process is the preferred implementation boundary.
The request would contain an operation name, named data or a checked
summary, units, sampling assumptions, resource limits and expected
source revision. It would not contain arbitrary R expressions. The
initial effect call uses the singleton local method on a verified CPDAG,
with \texttt{y.notparent=false}. The result preserves the multiset and
labels graph ambiguity separately from sampling uncertainty. Broader
PDAG constraints, optimal adjustment, joint interventions, PAG effects
and dependent observations require their own contracts. This restriction
makes the first experiment interpretable.

The package declares GPL version 2 or later, R at least 3.5.0, seventeen
imports and three \texttt{LinkingTo} dependencies. It requires
compilation; the selected build file links LAPACK, BLAS and Fortran
libraries. These are dependency declarations, not a resolved runtime
inventory. Separating R into a process helps isolate versions and
failures, but does not waive applicable source, notice or distribution
obligations. The actual dependency versions, build flags, numerical
libraries and third-party rights must be recorded for the artifact
selected for distribution \citep{exec13_pcalg_reuse_primary}.

The inspected source includes PC and IDA regression examples, including
comparisons across local/global methods, but some examples are
conditional on an extra-tests switch. Their presence supplies useful
candidate fixtures; it is not a passing test receipt. The inspected NEWS
begins with an older 2.7-8 section, and the separate ChangeLog begins in
2010. Neither provides a complete explanation of the selected release.
The build owner should therefore retain compatibility tests and a
revision-change review rather than treating the package date as a
support guarantee. The local inventory contains R 4.1.2 but no installed
\texttt{pcalg} or the inspected key graph/build packages; no package
benchmark is claimed here.

The reuse alternatives have different scientific and maintenance costs.

\begin{longtable}[]{@{}ll@{}}
\toprule
\begin{minipage}[b]{0.47\columnwidth}\raggedright
Choice\strut
\end{minipage} & \begin{minipage}[b]{0.47\columnwidth}\raggedright
Disposition for the proposed initial path\strut
\end{minipage}\tabularnewline
\midrule
\endhead
\begin{minipage}[t]{0.47\columnwidth}\raggedright
Direct native use\strut
\end{minipage} & \begin{minipage}[t]{0.47\columnwidth}\raggedright
Useful for a controlled R analysis, but a common product still needs
named graph, failure, unit and provenance semantics.\strut
\end{minipage}\tabularnewline
\begin{minipage}[t]{0.47\columnwidth}\raggedright
Process integration\strut
\end{minipage} & \begin{minipage}[t]{0.47\columnwidth}\raggedright
Preferred architecture after the pinned environment and strict boundary
pass their fixtures. It retains the native algorithm implementation
while isolating the transport.\strut
\end{minipage}\tabularnewline
\begin{minipage}[t]{0.47\columnwidth}\raggedright
Port the package\strut
\end{minipage} & \begin{minipage}[t]{0.47\columnwidth}\raggedright
Defer: reproducing orientation, numerical and multiset behavior would
create a new scientific maintenance obligation. The small exact
references are not a general port.\strut
\end{minipage}\tabularnewline
\begin{minipage}[t]{0.47\columnwidth}\raggedright
Reference only\strut
\end{minipage} & \begin{minipage}[t]{0.47\columnwidth}\raggedright
Appropriate now for documented capabilities and proposal comparisons;
insufficient for an executable workflow claim.\strut
\end{minipage}\tabularnewline
\begin{minipage}[t]{0.47\columnwidth}\raggedright
Defer qualification and distribution\strut
\end{minipage} & \begin{minipage}[t]{0.47\columnwidth}\raggedright
Required until environment, conversion and failure tests pass and the
resolved distribution is reviewed; local implementation can
proceed.\strut
\end{minipage}\tabularnewline
\begin{minipage}[t]{0.47\columnwidth}\raggedright
Reject permanently\strut
\end{minipage} & \begin{minipage}[t]{0.47\columnwidth}\raggedright
Unsupported by the present evidence; the identified hazards have
concrete, testable containment routes.\strut
\end{minipage}\tabularnewline
\bottomrule
\end{longtable}

The next useful experiment needs generated truth and a small pinned
environment, not expensive bank records. First build the isolated
dependency set and record its complete identity. Then exercise the named
chain, collider, undirected path, all column permutations, unit changes,
invalid covariance, orientation conflict, non-CPDAG and
exhausted-resource cases through the actual R process. Compare
well-conditioned CI calculations with the native test and verify the
deliberately stricter failure behavior separately. Only then run a
prospectively specified study of graph recovery, effect-set behavior and
complete-request cost. All refused and failed requests remain in its
denominator.

Installation time, process startup, data transfer, graph conversion,
effect enumeration, failed requests and revision maintenance belong in
the cost comparison. A faster local IDA call does not price the whole
workflow, and an interval spanning the candidate effects does not create
a sampling confidence guarantee. The numerical owner must justify the
strict callback; the integration owner must demonstrate representation
and failure preservation; the application owner must supply the
information clock and causal premises; and the distribution owner must
review the resolved artifact's obligations. Bank-specific validation
follows when a prototype result identifies a consequential unresolved
assumption and justifies the cost of testing it.

\label{exec13:pcalg-reuse-end}
